% Part: first-order-logic
% Chapter: syntax-and-semantics
% Section: terms-formulas

\documentclass[../../../include/open-logic-section]{subfiles}

\begin{document}

\olfileid{fol}{syn}{frm}

\olsection{ترم‌ها و \printtoken{P}{formula}}

پس از آنکه زبان مرتبهٔ اولی چون~$\Lang L$ داده شد، می‌توانیم عبارت‌هایی را
تعریف کنیم که از واژگان پایهٔ~$\Lang L$ ساخته شده‌اند. به‌ویژه،
\emph{ترم‌ها} و \emph{!!{formula}s} در زمرهٔ این عبارت‌ها هستند.

\begin{defn}[ترم‌ها]
\ollabel{defn:terms}
مجموعهٔ \emph{ترم‌های}~$\Trm[L]$ زبان~$\Lang L$ به‌صورت استقرایی چنین
تعریف می‌شود:
\begin{enumerate}
\item هر !!{variable} یک ترم است.
\item هر !!{constant} زبان~$\Lang L$ یک ترم است.
\item اگر $f$ یک !!{function} $n$جایگاهی و $t_1$، \dots، $t_n$ ترم باشند،
  آنگاه $\Atom{f}{t_1, \ldots, t_n}$ یک ترم است.
\tagitem{limitClause}{
  هیچ چیز دیگری ترم نیست.}{}
\end{enumerate}
ترمی که هیچ !!{variable}ی در بر نداشته باشد، \emph{ترم بسته} است.
\end{defn}

\begin{explain}
!!{constant}s در تعیین زبان و ترم‌های ما همچون رده‌ای جداگانه از نمادها
ظاهر می‌شوند، اما می‌شد آن‌ها را !!{function}sی صفرجایگاهی به شمار آورد.
در آن صورت می‌توانستیم بند دوم تعریف ترم‌ها را کنار بگذاریم. فقط باید
$\Atom{f}{t_1, \ldots, t_n}$ را، هرگاه $n = 0$ باشد، همان $f$ به‌تنهایی
بفهمیم.
\end{explain}

\begin{defn}[فرمول‌ها]
\ollabel{defn:formulas}
مجموعهٔ \emph{!!{formula}s}~$\Frm[L]$ زبان~$\Lang L$ به‌صورت استقرایی
چنین تعریف می‌شود:
\begin{enumerate}
\tagitem{prvFalse}{$\lfalse$ یک !!{formula} اتمی است.}{}

\tagitem{prvTrue}{$\ltrue$ یک !!{formula} اتمی است.}{}

\item اگر $R$ یک !!{predicate} $n$جایگاهیِ~$\Lang L$ و $t_1$، \dots،
  $t_n$ ترم‌های~$\Lang L$ باشند، آنگاه $\Atom{R}{t_1,\ldots, t_n}$ یک
  !!{formula} اتمی است.

\item اگر $t_1$ و $t_2$ ترم‌های~$\Lang L$ باشند، آنگاه
  $\Atom{\eq}{t_1, t_2}$ یک !!{formula} اتمی است.

\tagitem{prvNot}{اگر $!A$ یک !!a{formula} باشد، آنگاه $\lnot !A$ یک
  !!a{formula} است.}{}

\tagitem{prvAnd}{اگر $!A$ و $!B$ از !!{formula}s باشند، آنگاه $(!A \land
  !B)$ یک !!a{formula} است.}{}

\tagitem{prvOr}{اگر $!A$ و $!B$ از !!{formula}s باشند، آنگاه $(!A \lor !B)$
  یک !!a{formula} است.}{}

\tagitem{prvIf}{اگر $!A$ و $!B$ از !!{formula}s باشند، آنگاه $(!A \lif !B)$
  یک !!a{formula} است.}{}

\tagitem{prvIff}{اگر $!A$ و $!B$ از !!{formula}s باشند، آنگاه
  $(!A \liff !B)$ یک !!a{formula} است.}{}

\tagitem{prvAll}{اگر $!A$ یک !!a{formula} و $x$ یک !!a{variable} باشد،
  آنگاه $\lforall[x][!A]$ یک !!a{formula} است.}{}

\tagitem{prvEx}{اگر $!A$ یک !!a{formula} و $x$ یک !!a{variable} باشد،
  آنگاه $\lexists[x][!A]$ یک !!a{formula} است.}{}

\tagitem{limitClause}{هیچ چیز دیگری !!a{formula} نیست.}{}
\end{enumerate}
\end{defn}

\begin{explain}
تعریف مجموعهٔ ترم‌ها و تعریف مجموعهٔ !!{formula}s هر دو \emph{تعریف
استقرایی} هستند. در اصل، مجموعهٔ !!{formula}s را در بی‌نهایت مرحله
می‌سازیم. در مرحلهٔ آغازین، همهٔ فرمول‌های اتمی را فرمول اعلام می‌کنیم؛ این
مرحله با چند حالت نخست تعریف، یعنی حالت‌های
\iftag{prvTrue}{$\ltrue$، }{}%
\iftag{prvFalse}{$\lfalse$، }{}%
$\Atom{R}{t_1,\dots,t_n}$ و $\Atom{\eq}{t_1,t_2}$ متناظر است. پس
«!!{formula} اتمی» یعنی هر !!{formula}ی که یکی از این صورت‌ها را داشته باشد.

حالت‌های دیگر تعریف، قواعدی برای ساختن !!{formula}sی تازه از
!!{formula}sی پیش‌ترساخته به دست می‌دهند. در مرحلهٔ دوم می‌توانیم با آن
قواعد !!{formula}s را از !!{formula}sی اتمی بسازیم. در مرحلهٔ سوم،
فرمول‌های تازه را از فرمول‌های اتمی و فرمول‌های به‌دست‌آمده در مرحلهٔ دوم
می‌سازیم و همین‌طور ادامه می‌دهیم. اصطلاح «!!{formula}» بر هر چیزی اطلاق
می‌شود که سرانجام در یکی از این مراحل ساخته شود، و نه هیچ چیز دیگر.
\end{explain}

بنا بر قرارداد، $\eq$ را میان آرگومان‌هایش می‌نویسیم و پرانتزها را حذف
می‌کنیم: $\eq[t_1][t_2]$ اختصارِ $\Atom{\eq}{t_1,t_2}$ است. افزون بر این،
$\lnot \Atom{\eq}{t_1,t_2}$ را به‌اختصار $\eq/[t_1][t_2]$ می‌نویسیم.
هنگام نوشتن فرمول $(!B \ast !C)$ که از $!B$ و $!C$ با استفاده از رابط
دوجایگاهی~$\ast$ ساخته شده است، اغلب بیرونی‌ترین جفت پرانتز را حذف می‌کنیم
و فقط~$!B \ast !C$ می‌نویسیم.

\begin{intro}
برخی متون منطق لازم می‌دانند که !!{variable}~$x$ در~$!A$ رخ دهد تا
\iftag{prvEx}{$\lexists[x][!A]$ }{}%
\iftag{notprvEx,notprvAll}{}{و }%
\iftag{prvAll}{$\lforall[x][!A]$ }{}%
\iftag{notprvEx,notprvAll}{یک !!a{formula}}{!!{formula}s}
به شمار آیند. اگر این شرط را نگذارید هیچ مشکلی پیش نمی‌آید، و کار نیز
آسان‌تر می‌شود.
\end{intro}

\begin{tagblock}{defNot,defOr,defAnd,defIf,defIff,defTrue,defFalse,defEx,defAll}
\begin{defn}
فرمول‌های ساخته‌شده با عملگرهای تعریف‌شده را باید چنین فهمید:

\begin{tagenumerate}{defTrue,defFalse,defNot,defOr,defAnd,defIf,defIff,defEx,defAll}

\tagitem{defTrue}{$\ltrue$ اختصارِ
  \iftag{prvFalse}{$\lnot\lfalse$}{$(!A \lor \lnot !A)$ برای یک
    !!{formula} اتمی ثابت~$!A$} است.}{}

\tagitem{defFalse}{$\lfalse$ اختصارِ
  \iftag{prvTrue}{$\lnot\ltrue$}{$(!A \land \lnot !A)$ برای یک
    !!{formula} اتمی ثابت~$!A$} است.}{}

\tagitem{defNot}{$\lnot !A$ اختصارِ $!A \lif \lfalse$ است.}{}

\tagitem{defOr}{$!A \lor !B$ اختصارِ
  \iftag{prvAnd}{$\lnot(\lnot !A \land \lnot !B)$}{$\lnot !A \lif
    !B$} است.}{}

\tagitem{defAnd}{$!A \land !B$ اختصارِ
  \iftag{prvOr}{$\lnot(\lnot !A \lor \lnot !B)$}{$\lnot (!A \lif
    \lnot !B)$} است.}{}

\tagitem{defIf}{$!A \lif !B$ اختصارِ
  \iftag{prvOr}{$(\lnot !A \lor !B)$}{$\lnot (!A \land \lnot !B)$} است.}{}

\tagitem{defIff}{$!A \liff !B$ اختصارِ $(!A \lif !B) \land (!B
  \lif !A)$ است.}{}

\tagitem{defAll}{$\lforall[x][!A]$ اختصارِ $\lnot\lexists[x][\lnot !A]$ است.}{}

\tagitem{defEx}{$\lexists[x][!A]$ اختصارِ $\lnot\lforall[x][\lnot !A]$ است.}{}
\end{tagenumerate}
\end{defn}
\end{tagblock}

اگر در زبانی برای کاربردی خاص کار کنیم، اغلب !!{predicate}s و !!{function}sی
دوجایگاهی را میان ترم‌های مربوط می‌نویسیم؛ برای نمونه، $t_1 < t_2$ و
$(t_1 + t_2)$ در زبان حساب، و $t_1 \in t_2$ در زبان نظریهٔ مجموعه‌ها.
تابع جانشین در زبان حساب را حتی بنا بر قرارداد \emph{پس از} آرگومانش
می‌نویسند:~$t'$. بااین‌حال، صورت رسمی این عبارت‌ها به‌ترتیب صرفاً
اختصارهای $\Atom{\Obj A^2_0}{t_1, t_2}$، $\Obj f^2_0(t_1, t_2)$،
$\Atom{\Obj A^2_0}{t_1, t_2}$ و $\Obj f^1_0(t)$ است.

\begin{defn}[همانی نحوی]
نماد $\ident$ همانی نحوی میان رشته‌های نمادها را بیان می‌کند؛ یعنی
$!A \ident !B$ اگر و تنها اگر $!A$ و $!B$ رشته‌هایی از نمادها باشند که
طول یکسان دارند و در هر جایگاه نماد یکسانی دارند.
\end{defn}

در دو سوی نماد $\ident$ می‌توان رشته‌هایی را قرار داد که با الحاق به دست
آمده‌اند. برای نمونه، $!A \ident (!B \lor !C)$ یعنی رشتهٔ نمادهای~$!A$
همان رشته‌ای است که با الحاق یک پرانتز باز، رشتهٔ $!B$، نماد $\lor$، رشتهٔ
$!C$ و یک پرانتز بسته، به همین ترتیب، به دست می‌آید. اگر چنین باشد، می‌دانیم
نخستین نماد~$!A$ یک پرانتز باز است، $!A$ رشتهٔ $!B$ را به‌عنوان زیررشته‌ای
در بر دارد که از نماد دوم آغاز می‌شود، پس از آن زیررشته نماد~$\lor$ می‌آید،
و به همین ترتیب.

چون ترم‌ها و !!{formula}s با تعریف‌های استقرایی از عناصر پایه ساخته
می‌شوند، می‌توانیم برای اثبات ویژگی‌های آن‌ها از اصول استقرای زیر بهره
بگیریم.

\begin{lem}[\emph{اصل استقرا بر ترم‌ها}]
    \ollabel{lem:trmind}
    بگذارید $\Lang L$ زبانی مرتبهٔ اول باشد.
    اگر ویژگی‌ای چون~$P$ چنان باشد که
    %
    \begin{enumerate}
        \item برای هر !!{variable}~$v$ برقرار باشد،
        %
        \item برای هر !!{constant}~$a$ زبان~$\Lang L$ برقرار باشد، و
        %
        \item هرگاه برای $t_1$،~\dots، $t_n$ برقرار باشد و $f$ یک
            !!{function} $n$جایگاهیِ~$\Lang L$ باشد، برای
            $f(t_1,\dotsc,t_n)$ نیز برقرار باشد
    \end{enumerate}
    (با فرض اینکه $t_1$،~\dots، $t_n$ ترم‌های~$\Lang{L}$ باشند)،
    آنگاه $P$ برای هر ترم در~$\Trm[L]$ برقرار است.
\end{lem}

\begin{prob}
    \olref[fol][syn][frm]{lem:trmind} را اثبات کنید.
\end{prob}

\begin{lem}[\emph{اصل استقرا بر !!{formula}s}]
    \ollabel{thm:frmind}
    بگذارید $\Lang L$ زبانی مرتبهٔ اول باشد.
    اگر ویژگی‌ای چون~$P$ برای همهٔ !!{formula}sی اتمی برقرار باشد و
    چنان باشد که
    %
    \begin{enumerate}
        \tagitem{prvNot}{هرگاه برای~$!A$ برقرار باشد، برای $\lnot !A$
            نیز برقرار باشد؛}{}
        \tagitem{prvAnd}{هرگاه برای $!A$ و~$!B$ برقرار باشد، برای
            $(!A \land !B)$ نیز برقرار باشد؛}{}
        \tagitem{prvOr}{هرگاه برای $!A$ و~$!B$ برقرار باشد، برای
            $(!A \lor !B)$ نیز برقرار باشد؛}{}
        \tagitem{prvIf}{هرگاه برای $!A$ و~$!B$ برقرار باشد، برای
            $(!A \lif !B)$ نیز برقرار باشد؛}{}
        \tagitem{prvIff}{هرگاه برای $!A$ و~$!B$ برقرار باشد، برای
            $(!A \liff !B)$ نیز برقرار باشد؛}{}
        \tagitem{prvEx}{هرگاه برای~$!A$ برقرار باشد، برای
            $\lexists[x][!A]$ نیز برقرار باشد؛}{}
        \tagitem{prvAll}{هرگاه برای~$!A$ برقرار باشد، برای
            $\lforall[x][!A]$ نیز برقرار باشد؛}{}
  \end{enumerate}
  (با فرض اینکه $!A$ و $!B$ از !!{formula}sی~$\Lang{L}$ باشند)،
  آنگاه $P$ برای همهٔ فرمول‌های~$\Frm[L]$ برقرار است.
\end{lem}

\end{document}
